\documentclass{minmal}
\usepackage{graphicx} % Required for inserting images
\usepackage{tikz}
\usepackage[left=2.00cm, right=1.00cm]{geometry}
\usepackage{xcolor} % before tikz or tkz-euclide if necessary
\usepackage{amsmath}
\usepackage{siunitx}
\usepackage{chemfig}
\usetikzlibrary{calc}
\usepackage{enumitem}
\begin{document}
\begin{center}
    \large{\textbf{Prøve i kjemisk teknologi:}}
\end{center}
\textbf{Oppg. 1:}\\
\liste[a]{
	\item Hvilke tre bergarter har vi?
	\item Gi eksempel på 3 ulike mineraler.  Hvilke metaller kan vi fremstille fra dem?
	\item Hvilke egenskaper er vanlige for metallene?
	\item Hva menes med et lett metall og et tungmetall (hint:  ta utgangspunkt i metallets tetthet).
}
\textbf{Oppg. 2:}\\
    \liste[*]{
        \item Hvor mange protoner, og hvor mange elektroner har karbon?  
        \item Tegn elektronprikkstrukturen for karbon.
        \item Skriv formelen molekylene metan, propan og butan.
        \item Hva har kull hovedsakelig blitt dannet av?
    }
\textbf{Oppg 3:}\nl[1]
\liste[a]{
\item Bildet underviser en enkel fremstilling av et vannkraftverk:\\
\begin{figure}[h]
    \centering
    \includegraphics[width=0.5\linewidth]{vannkraftverk.png}
    \label{fig:placeholder}
\end{figure}\\
Gi en enkel forklaring på hvordan vannkraftverket fungerer.
\item Når hydrogen og oksygen forbrenner i en hydrogendrevet bil skje følgende reaksjon:

		$$2H_{2}+O_{2}=2H_{2}O(g)+energi$$
	Hva består eksosen til hydrogenbilen av?
\item Hva brukes energien som dannes til?

\item Hva får Borregård produsert fra tømmerstokken?
}
\textbf{Oppg. 4:}\nl[1]
\liste[a]{
\item Hvilke prinsipper har vi for oppdeling av faste stoffer?
\item Forklar virkemåten til en kjeftknuser:\\
\begin{figure}[h]
    \centering
    \includegraphics[width=0.5\linewidth]{kjeftknuser.png}
    \label{0}
\end{figure}
\newpage
\item Forklar virkemåten til en hammermølle:\nl[1]
\textbf{}
\begin{figure}[h]
    \centering
    \includegraphics[width=0.5\linewidth]{hammermølle.png}
    \label{fig:placeholder}
\end{figure}
\item Forklar virkemåten til en kulemølle:
 \begin{figure}[h]
     \centering
     \includegraphics[width=0.5\linewidth]{kulemølle.png}
     \label{fig:placeholder}
 \end{figure}

Forklar også hvilken hastighet kulemøllen bør ha.
}
\newpage
\textbf{Oppg. 5:}\nl[1]
\liste[a]{
    \item Forklar virkemåten til en granuler tallerken:\nl[1]
    \begin{figure}[h]
        \centering
        \includegraphics[width=0.75\linewidth]{granulertallerken.png}
        \label{fig:placeholder}
    \end{figure}

 


	\item Hva menes med lavine-effekten?
}	
\end{document}